% Pursuing Stacks (cotes 134-2 to 134-8): archive cote 134-6, site page 19 (PDF page 20), typescript page 338.
% typescript. Draft compared: ps5.tex:2322-2388.
% AI-transcribed by Claude Opus 5.5 (claude-opus-5-5, adaptive thinking, effort high), Message Batches API, batch ps3 (msgbatch_01XdFmuVVBAcgc7EFxdtCqsj), 2026-11-12; not yet checked by eye.
\scanpage{134-6/19}
\mspage{338}
replace it by a suitable substitute,  such as the existence,  for any
two objects a and b in A,  of a cellular subdivision of $\hand{|}i(a)\hand{|}\times\hand{|}i(b)\hand{|}$ ,
made up with cells of the type $\hand{|}i(c)\hand{|}$,  and inducing on the latter the
given cellular structure of $\hand{|}i(c)\hand{|}$ .  The problem is now (besides get-%
ting or not an abelianizator $L.$ )  \uline{whether A is at any rate a weak test
category}  (in view of the example of $\hand{\Delta}^{f}$ ,  we can't expect\inserted{now} of course
\xout{\illegible{1}} to get an actual test category).  Maybe I'll come back to this la-%
ter,  when writing down a proof for $\hand{\Delta}^{f}$  being a weak test category,  if
a more general result along these lines should come out alongside.

\uline{Remarks}.  \deleted{\unsure{1)}}  I feel the canonical chain complex $L.$  in  $\mathrm{Add}(A^{o})$  con-%
structed in this section,  under suitable assumptions on the small cate-%
gory A,  merits a name of it's own.  We may call it the \uline{standard
abelianizator}  of A  -  but this is reasonable only if it turns out that
in al\over{\unsure{d}}{l} cases when it can be constructed,  it is an abelianizator indeed.
Another convenient name may be the \uline{Dold-Puppe} \uline{Chain complex} ,  as \xout{the}
in the three standard cases\hand{,} the standard Dold-Puppe construction of the
``normalization''  of an abelian complex\inserted{(ss say)} can be viewed as being performed
in the ``universal'' case,  namely \xout{xx}  $A^{o} \hand{\longrightarrow} \mathrm{Add}(A^{o})$,  and the correspon-%
ding ``full'' chain complex,  namely  $L.$  -  with this \xout{slight} grain of
salt though that we have still to enlarge  $\mathrm{Add}(A^{o})$ slightly,  so as to
make it stable under taking direct summands corresponding to projectors.
But then it occurs to me that the name of Dold-Puppe chain complex is
much more suitable for the \uline{result of normalization}  applied to $L.$  ,
\margin{$\uparrow$}which (if I got it right)  is the ``new'' complex discovered by Dold-Pup-%
pe,  together with the inverse construction,  whereas $L.$  had already\inserted{been} \deleted{be}
known for ages (even if not under it's universal disguise ...).

\margin{$\downarrow$}\hand{94.} \deleted{\illegible{1}}  In \deleted{this}\inserted{the last} section,  as in the two preceding days,  our emphasis
\xout{has been on} with abelianization of homotopy types has been to look at it
in terms of more or less arbitrary test categories and the corresponding
elementary modelizers,  are even in terms of arbitrary small categories.
This has caused as spinning a kind of dream for a while\hand{,}with the White-%
head and Dold Puppe theorems  and generalized bou\inserted{n}dary maps as our main
thread.   Now this reminds me of a rather different line of thoughts
tied up with abelianization,  quite independently of playing around
with variable modelizers  -  a question which has been intriguing me
for a very long time now,  ever since  I  got acquainted a little with
the very notion of homotopy types,  and the corresponding homology and

\margin{19}
